\chapter{Convergence of a min--max sequence to a minimal hypersurface}\label{ch:convergence-of-min-max}

This chapter completes the proof of the main theorem by showing that the min-max sequence constructed in the previous chapters converges to a smooth free-boundary minimal hypersurface. The argument combines three key ingredients: the pull-tight procedure adapted to the complete setting, the almost minimizing property in relative annuli, and the free-boundary regularity theory of Li--Zhou \cite{LZ17}.

Throughout this chapter we work in the doubled manifold $\widetilde M$ and use vector fields that are tangent to $\partial M$ in a collar neighbourhood. Varifold convergence is always understood in $\widetilde M$, so "stationary" means stationary with free boundary in $M$.

Starting from a good sweepout of the good set constructed in Chapter~\ref{ch:nested-sweepout}, we perform a pull--tight deformation on an exhaustion of $\widetilde M$, select slices that are almost minimizing in small relative annuli, and pass to the limit. The free boundary regularity theory of Li--Zhou~\cite{LZ17}, together with the almost minimizing property, then yields a smooth almost properly embedded free boundary minimal hypersurface.

Our treatment of non-excessive optimal sweepouts and the auxiliary quantities $\mathfrak{h}_{\mathrm{nm}}$ and $\Snm$ follows the framework developed by Chodosh--Liokumovich--Spolaor; in particular we adopt the notation and conventions from \cite{CLS22} throughout this chapter.

\section{Manifolds with sublinear growth}\label{sec:sublinear-width}
The following lemma adapts \cite[§8.1]{CL20} to the relative setting.

\begin{lemma}\label{lem:good-set-sublinear}
Let $(M^{n+1},\partial M)$ be complete with $3\le n+1\le 7$ and sublinear volume growth, i.e.
\[
\liminf_{R\to\infty}\frac{\haus^{n+1}\bigl(B_R(x)\cap M\bigr)}{R}=0
\qquad\text{for some }x\in M.
\]
Then there exists a good set $U\subset M$ with $0<W_g(U)<\infty$.
\end{lemma}

\begin{proof}
Without loss of generality, we may assume $x \in \mathrm{int}(M)$ (if $x \in \partial M$, replace it by a nearby interior point). Choose $r>0$ so that $B_r(x)\Subset \mathrm{int}(M)$. The local isoperimetric inequality provides a constant
\[
C_I:=\inf\Bigl\{\haus^n(\partial E\cap B_r(x)):\ E\subset B_r(x),\
\haus^{n+1}(E)=\tfrac12\,\haus^{n+1}(B_r(x))\Bigr\}>0.
\]
The coarea formula yields $\frac{\mathrm{d}}{\mathrm{d}R}\haus^{n+1}(B_R(x)\cap M)=\haus^n(\partial B_R(x)\cap M)$ for almost every $R>0$.  The sublinear assumption therefore allows us to pick a regular value $R>r$ with
\[
\haus^n(\partial B_R(x)\cap M)\le \frac{C_I}{100}.
\]
Set $U:=B_R(x)\cap M$.  The relative boundary of $U$ equals $\partial B_R(x)\cap M$, so $\haus^{n}(\partial_{\mathrm{rel}}U)\le C_I/100$.

Let $\{\Gamma_t\}$ be any sweepout of $U$ with corresponding regions $\{\Omega_t\}$.  The function $h(t):=\haus^{n+1}(\Omega_t\cap B_r(x))$ is continuous and takes the value $\tfrac12\haus^{n+1}(B_r(x))$, hence one of the slices bisects volume in $B_r(x)$.  For that slice, $\haus^{n}(\Gamma_t\cap M)\ge C_I$.  Consequently $W_\partial(U)\ge C_I$, while $\haus^{n}(\partial_{\mathrm{rel}}U)\le C_I/100$, and $U$ is a good set.

To bound $W_g(U)$ from above, smooth the distance function on $\widetilde M$ to a Morse function $\tilde d_x$ with $|\nabla \tilde d_x|\le 1$ and tangent to $\partial M$ in a collar.  The level sets $\Gamma_s=\{\tilde d_x=s\}$ give a nested sweepout of $U$ whose area is controlled by $\sup_{s\in[0,R]}\haus^n(\partial B_s(x)\cap M)$.  Adjusting the endpoints inside a fixed collar of $\partial U$ produces a good sweepout, showing that $W_g(U)<\infty$.  Finally, the bisecting argument above implies $W_g(U)\ge C_I$.
\end{proof}

\section{Pulled--tight sequences and min--max varifolds}
Fix a good set $U\subset M$.  Proposition~\ref{prop:nested-sweepout} produces good sweepouts with maximal area arbitrarily close to $W_g(U)$.  To prevent loss of mass at infinity we perform the pull--tight deformation on a bounded exhaustion of $\widetilde M$.

Let $U_0\Subset U_1\Subset\cdots\Subset \widetilde M$ be bounded relative open sets with $M=\bigcup_i U_i$ and $U\Subset U_0$.  For each $i$ choose $W_i$ with $U_i\Subset W_i\Subset U_{i+1}$.  For a bounded relative open set $\Omega\subset M$ denote by $\mathcal{V}_\Omega$ the set of varifolds supported in $Gr_n(\Omega)$ with mass $\le 2W_g(U)$, and let $\mathcal{V}_{\Omega,\mathrm{st}}$ be the subset stationary with free boundary in $\Omega$.

\begin{lemma}[Extension of the pull--tight map]\label{lem:extend-pull-tight}
For every $i$ there exist a continuous map $\Phi_{U_i,W_i}:\mathcal{V}(M)\to\mathcal{V}(M)$ and decreasing sequences $\{\tau_k(U_i)\}$, $\{\varepsilon_k(U_i)\}$ converging to $0$ with the following properties for all $V\in\mathcal{V}(M)$:
\begin{enumerate}
    \item $\|\Phi_{U_i,W_i}(V)\|(M)\le \|V\|(M)$;
    \item if $\|V\|(\overline{W}_i)\le 5\,\haus^n(\partial_{\mathrm{rel}}U)$ then $\Phi_{U_i,W_i}(V)=V$;
    \item if $\|V\|(\overline{W}_i)\ge 9\,\haus^n(\partial_{\mathrm{rel}}U)$ and
    \[
    \mathbf{F}_{U_i}\bigl(V\llcorner Gr_n(U_i),\mathcal{V}_{U_i,\mathrm{st}}\bigr)\in\bigl[2^{-(k+1)},2^{-k}\bigr],
    \]
    then $\|\Phi_{U_i,W_i}(V)\|(M)\le \|V\|(M)-\varepsilon_k(U_i)$ and $\mathbf{F}\bigl(\Phi_{U_i,W_i}(V),V\bigr)\le \tau_k(U_i)$.
\end{enumerate}
\end{lemma}

\begin{proof}
The argument is the same as in \cite[Lemma~8.8]{CL20}, adapted to the free-boundary setting following Li--Zhou~\cite[Section~4]{LZ17}.  Admissible vector fields are taken tangent to $\partial M$ by averaging over the reflection in the doubled manifold.  The compactness of the sets
\[
\Bigl\{V\in\mathcal{V}_{U_i}: \mathbf{F}_{U_i}(V,\mathcal{V}_{U_i,\mathrm{st}})\in\bigl[2^{-(k+1)},2^{-k}\bigr]\Bigr\}
\]
allows us to choose finitely many vector fields whose flows decrease area by a fixed amount; composing them with a partition of unity gives the desired map.
\end{proof}

Let $\{\Gamma_t^i\}$ be a sequence of good sweepouts with $\sup_t \haus^n(\Gamma_t^i\cap M)\downarrow W_g(U)$.  Set
\[
\mathbf{P}_i:=\Phi_{U_0,W_0}\circ\Phi_{U_1,W_1}\circ\cdots\circ\Phi_{U_i,W_i}.
\]

\begin{proposition}[Pulled--tight minimizing sequence]\label{prop:pt-minimizing}
After passing to a subsequence, $\{\mathbf{P}_i(\Gamma_t^i)\}$ is a minimizing sequence of good sweepouts with
\[
\lim_{i\to\infty}\ \sup_{t\in[0,1]}\haus^n\bigl(\mathbf{P}_i(\Gamma_t^i)\cap M\bigr)=W_g(U),
\]
and any min--max sequence $\mathbf{P}_i(\Gamma_{t_i}^i)$ admits a subsequence converging to a varifold stationary with free boundary in $M$.
\end{proposition}

\begin{proof}
Property (3) in Lemma~\ref{lem:extend-pull-tight} guarantees that if the restriction of a slice to some $U_j$ stays a definite $\mathbf{F}_{U_j}$--distance from $\mathcal{V}_{U_j,\mathrm{st}}$, then the composed map decreases its mass by at least $\varepsilon_k(U_j)$.  Selecting a diagonal subsequence therefore produces slices whose restrictions to every $U_j$ converge to the stationary set, hence the limit varifold is stationary with free boundary.
\end{proof}

\section{Almost minimizing slices in annuli}\label{sec:am-annuli}
Let $\mathcal{A}(r,U)$ denote the family of balls and annuli of radius $\le r$ intersecting $U$, and set
\begin{align*}
\mathrm{CO}(\mathcal{A}(r,U)):=\Bigl\{(A_1,A_2):A_i\in\mathcal{A}(r,U),\ \\
\inf_{x\in A_1,y\in A_2}\operatorname{dist}(x,y)\ge 4\min\{\operatorname{diam}A_1,\operatorname{diam}A_2\}\Bigr\}.
\end{align*}
We make repeated use of the continuous `almost minimizing' notion in Definition~\ref{def:AM-rel} from Appendix~\ref{appendix:AM}.

\begin{lemma}[Freezing and dynamic competitor]\label{lem:freeze-dynamic-rel}
Let $U\Subset U'\Subset M$ be relative open sets and $\{\partial\Xi_t\}$ a good sweepout.  Suppose there exists $t_0\in[0,1]$ and a one-parameter family $\{\partial\Omega_s\}_{s\in[0,1]}$ satisfying the sweepout regularity conditions (s1)--(s4), (sw1)--(sw3) with $\Omega_0=\Xi_{t_0}$, $\Omega_s\setminus U'=\Xi_{t_0}\setminus U'$ for all $s$, and
\begin{enumerate}[label=(\roman*)]
    \item $\haus^n(\partial\Omega_s\cap M)\le \haus^n(\partial\Xi_{t_0}\cap M)+\varepsilon/8$ for all $s\in[0,1]$;
    \item $\haus^n(\partial\Omega_1\cap M)\le \haus^n(\partial\Xi_{t_0}\cap M)-\varepsilon$.
\end{enumerate}
Then there exists $\eta>0$ such that for every $a,b,a',b'$ with
\[
t_0-\eta\le b<b'<a'<a\le t_0+\eta,
\]
there is a homotopic competitor sweepout $\{\partial\Xi_t'\}$ satisfying:
\begin{enumerate}[label=(\alph*)]
    \item $\Xi_t'=\Xi_t$ for $t\in[0,b]\cup[a,1]$ and $\Xi_t'\setminus U'=\Xi_t\setminus U'$ for $t\in(b,a)$;
    \item $\haus^n(\partial\Xi_t'\cap M)\le \haus^n(\partial\Xi_t\cap M)+\varepsilon/4$ for all $t$;
    \item $\haus^n(\partial\Xi_t'\cap M)\le \haus^n(\partial\Xi_t\cap M)-\varepsilon/2$ for $t\in[a',b']$.
\end{enumerate}
\end{lemma}

\begin{proof}
The proof adapts the freezing and dynamic competitor construction of \cite[Lemma~3.1]{DLT13} to the free-boundary setting. We work throughout in the doubled manifold $\widetilde M$, which allows all deformations to be generated by vector fields tangent to $\partial M$.

\noindent\textbf{Step 1: Freezing.}
Choose relative open sets $A$ and $B$ such that
\begin{itemize}
    \item $U\Subset A\Subset B\Subset U'$;
    \item $\partial\Xi_{t_0}\cap C$ is smooth, where $C:=B\setminus\overline A$.
\end{itemize}
This is possible since $\partial\Xi_{t_0}$ has only finitely many singular points. Next, fix smooth functions $\varphi_A\in C_c^\infty(B)$ and $\varphi_B\in C_c^\infty(M\setminus\overline A)$ with $\varphi_A+\varphi_B\equiv 1$.

By the smooth convergence of $\Xi_t$ to $\Xi_{t_0}$ as $t\to t_0$, we can fix $\eta>0$ and $C'\Subset C$ such that for every $t\in(t_0-\eta,t_0+\eta)$:
\begin{itemize}
    \item $\partial\Xi_t\cap C$ is the graph of a function $g_t$ over $\partial\Xi_{t_0}\cap C$ (using Fermi coordinates);
    \item $\Xi_t\cap(C\setminus C')=\Xi_{t_0}\cap(C\setminus C')$;
    \item $\Xi_t\cap C'=\{(z,\sigma):\sigma<g_t(z)\}\cap C'$.
\end{itemize}
Note that $g_{t_0}\equiv 0$ by construction. Define
\[
g_{t,s,\tau}:=\varphi_B g_t+\varphi_A\bigl((1-s)g_t+sg_\tau\bigr)
\quad\text{for }t,\tau\in(t_0-\eta,t_0+\eta),\;s\in[0,1].
\]
Since $g_t\to g_{t_0}$ smoothly as $t\to t_0$, we can choose $\eta$ small enough so that
\[
\sup_{s,\tau}\|g_{t,s,\tau}-g_t\|_{C^1}
\]
is arbitrarily small. In particular, if $\Gamma_{t,s,\tau}$ denotes the graph of $g_{t,s,\tau}$, we can ensure
\[
\max_s\haus^n(\Gamma_{t,s,\tau}\cap C)\le \haus^n(\partial\Xi_t\cap C)+\varepsilon/16.
\]

Given $t_0-\eta<b<b'<a'<a<t_0+\eta$, choose $b''\in(b, b')$ and $a''\in(a',a)$, and fix:
\begin{itemize}
    \item a smooth function $\psi:[b,a]\to[0,1]$ which is identically $0$ near $b$ and $a$, and identically $1$ on $[b'',a'']$;
    \item a smooth function $\gamma:[b,a]\to[t_0-\eta,t_0+\eta]$ which equals the identity near $b$ and $a$, and is constantly $t_0$ on $[b'',a'']$.
\end{itemize}

Define $\{\Delta_t\}$ by:
\begin{itemize}
    \item $\Delta_t=\Xi_t$ for $t\notin[b,a]$;
    \item $\Delta_t\setminus\overline B=\Xi_t\setminus\overline B$ for all $t$;
    \item $\Delta_t\cap A=\Xi_{\gamma(t)}\cap A$ for $t\in[b,a]$;
    \item $\Delta_t\cap(C\setminus C')=\Xi_{t_0}\cap(C\setminus C')$;
    \item $\Delta_t\cap C'=\{(z,\sigma):\sigma<g_{t,\psi(t),\gamma(t)}(z)\}$ for $t\in[b,a]$.
\end{itemize}
Then $\{\partial\Delta_t\}$ is a sweepout homotopic to $\{\partial\Xi_t\}$ with:
\begin{itemize}
    \item $\Delta_t=\Xi_t$ outside $[b,a]$ and $\Delta_t\setminus U'=\Xi_t\setminus U'$ for all $t$;
    \item $\Delta_t\cap U=\Xi_{t_0}\cap U$ for $t\in[b'',a'']$ (i.e., \textbf{frozen} in this interval);
    \item $\haus^n(\partial\Delta_t\cap C)\le \haus^n(\partial\Xi_t\cap C)+\varepsilon/16$ for $t\in[b,a]$.
\end{itemize}

\noindent\textbf{Step 2: Dynamic competitor.}
Fix a smooth function $\chi:[b'',a'']\to[0,1]$ which is identically $0$ near $b''$ and $a''$, and identically $1$ on $[b',a']$. Set:
\begin{itemize}
    \item $\Xi_t'=\Delta_t$ for $t\notin[b'',a'']$;
    \item $\Xi_t'\setminus A=\Delta_t\setminus A$ for $t\in[b'',a'']$;
    \item $\Xi_t'\cap A=\Omega_{\chi(t)}\cap A$ for $t\in[b'',a'']$.
\end{itemize}
Then $\{\partial\Xi_t'\}$ is a sweepout homotopic to $\{\partial\Xi_t\}$. Property (a) is immediate. For $t\notin[b,a]$, $\Xi_t'=\Xi_t$, so property (b) holds trivially. For $t\in[b,a]$, we have $\Xi_t'=\Xi_t$ on $M\setminus B$ and $\Xi_t'=\Delta_t$ on $C$, giving
\begin{equation}\label{eq:freeze-area-estimate}
\haus^n(\partial\Xi_t'\cap M)-\haus^n(\partial\Xi_t\cap M)
\le \frac{\varepsilon}{16}+\bigl[\haus^n(\partial\Xi_t'\cap A)-\haus^n(\partial\Xi_t\cap A)\bigr].
\end{equation}

We estimate the change in $A$ in three cases:

\begin{enumerate}[label=(\roman*)]
\item For $t\in[b,b'']\cup[a'',a]$: Here $\Xi_t'\cap A=\Xi_{\gamma(t)}\cap A$. Choosing $\eta$ small, we have
\[
\bigl|\haus^n(\partial\Xi_s\cap A)-\haus^n(\partial\Xi_\sigma\cap A)\bigr|\le \varepsilon/16
\quad\text{for all }s,\sigma\in(t_0-\eta,t_0+\eta).
\]
Thus $\haus^n(\partial\Xi_t'\cap M)\le \haus^n(\partial\Xi_t\cap M)+\varepsilon/8$.

\item For $t\in[b'',b']\cup[a',a'']$: Here $\partial\Xi_t'\cap A=\partial\Omega_{\chi(t)}\cap A$. Using assumption (i) and the estimate for $\partial\Xi_t\cap A$:
\[
\haus^n(\partial\Xi_t'\cap M)-\haus^n(\partial\Xi_t\cap M)
\le \frac{\varepsilon}{16}+\frac{\varepsilon}{16}+\frac{\varepsilon}{8}=\frac{\varepsilon}{4}.
\]

\item For $t\in[b',a']$: Here $\Xi_t'\cap A=\Omega_1\cap A$. Using assumption (ii):
\[
\haus^n(\partial\Xi_t'\cap M)-\haus^n(\partial\Xi_t\cap M)
\le \frac{\varepsilon}{16}+\bigl[\haus^n(\partial\Omega_1\cap A)-\haus^n(\partial\Xi_{t_0}\cap A)\bigr]
<-\frac{\varepsilon}{2}.
\]
\end{enumerate}
This establishes (b) and (c).
\end{proof}

\begin{proposition}[Almost minimizing selection]\label{prop:AM-selection}
Let $\{\Gamma_t^i\}$ be a minimizing sequence as in Proposition~\ref{prop:pt-minimizing} and assume $\sup_{t}\haus^n(\Gamma_t^i\cap M)<W_g(U)+1/(8i)$.  For every $r>0$ there exists $i_0(r)$ such that for $i\ge i_0(r)$ one can choose $t_i\in[0,1]$ with
\begin{enumerate}
    \item $\Gamma_i:=\mathbf{P}_i(\Gamma_{t_i}^i)$ is $1/i$--almost minimizing in every pair $(A_1,A_2)\in\mathrm{CO}(\mathcal{A}(r,U))$;
    \item $\haus^n(\Gamma_i\cap M)\ge W_g(U)-1/i$;
    \item $\haus^n(\Gamma_i\cap N_r(U))\ge \varepsilon(U)/2$, where $\varepsilon(U)$ is provided by Proposition~\ref{prop:non-escape}.
\end{enumerate}
\end{proposition}

\begin{proof}
We argue by contradiction. Assume $i$ is sufficiently large so that $1/i<\varepsilon(U)/2$. Let
\[
A_i:=\bigl\{t\in[0,1]:\haus^n(\Gamma_t^i\cap M)\ge W_g(U)-1/i\bigr\}
\]
and
\[
B_i(U,r):=\bigl\{t\in[0,1]:\haus^n(\Gamma_t^i\cap N_r(U))\ge \varepsilon(U)/2\bigr\}.
\]
Define $K_i(U,r):=A_i\cap B_i(U,r)$. This is a compact set since $A_i$ and $B_i(U,r)$ are both closed. By Proposition~\ref{prop:non-escape}, $K_i(U,r)$ is non-empty.

Suppose the proposition is false. Then there exists a sequence $i_k\to\infty$ such that $\Gamma_t^{i_k}$ is not $1/i_k$--almost minimizing in some pair $(U_t^1,U_t^2)\in\mathrm{CO}(\mathcal{A}(r,U))$ for every $t\in K_{i_k}(U,r)$. To simplify notation, we drop the subscript $i_k$ and work with a fixed large index.

We will modify the family $\{\Gamma_t\}$ on an open set containing $K:=K(U,r)\subset[0,1]$ to produce a new family $\{\Gamma_t'\}$ with $\haus^n(\Gamma_t'\cap M)<W_g(U)$ for all $t$ where $\haus^n(\Gamma_t'\cap N_r(U))>\varepsilon(U)/2$. This contradicts Proposition~\ref{prop:non-escape}.

By Lemma~\ref{lem:freeze-dynamic-rel} and the covering refinement argument of \cite[page~13]{DLT13}, we can choose:
\begin{itemize}
    \item a covering $\{J_j=(b_j,a_j)\}$ of $K$ such that each point of $K$ is contained in at most two intervals $J_j$;
    \item a collection of sets $\{U_j\}\subset\mathcal{A}(r,U)$ such that if $\overline{J_j}\cap\overline{J_{j'}}\neq\varnothing$, then $\inf_{x\in U_j,y\in U_{j'}}\operatorname{dist}(x,y)>0$;
    \item $\delta>0$ such that $\{(b_j+\delta,a_j-\delta)\}$ still covers $K$;
    \item for each $j$, a family $\{\Omega_{j,t}\}$ obtained from Lemma~\ref{lem:freeze-dynamic-rel} with:
    \begin{enumerate}[label=(\alph*)]
        \item $\Omega_{j,t}=\Omega_t$ for $t\notin J_j$ and $\Omega_{j,t}\setminus U_j=\Omega_t\setminus U_j$ for all $t$;
        \item $\haus^n(\partial\Omega_{j,t}\cap M)\le \haus^n(\partial\Omega_t\cap M)+1/(4i)$ for all $t$;
        \item $\haus^n(\partial\Omega_{j,t}\cap M)\le \haus^n(\partial\Omega_t\cap M)-1/(2i)$ for $t\in(b_j+\delta,a_j-\delta)$.
    \end{enumerate}
\end{itemize}

Define a new sweepout $\{\partial\Omega_t'\}$ of $U$ by:
\begin{itemize}
    \item $\Omega_t'=\Omega_t$ if $t\notin\bigcup_j(b_j,a_j)$;
    \item $\Omega_t'=\Omega_{j,t}$ if $t$ is contained in a single $J_j$;
    \item $\Omega_t'=\bigl[\Omega_t\setminus(U_j\cup U_{j+1})\bigr]\cup[\Omega_{j,t}\cap U_j]\cup[\Omega_{j+1,t}\cap U_{j+1}]$ if $t\in J_j\cap J_{j+1}$.
\end{itemize}

\noindent\textbf{Claim:} If $\haus^n(\partial\Omega_t'\cap N_r(U))\ge \varepsilon(U)/2$, then $\haus^n(\partial\Omega_t'\cap M)<W_g(U)$.

The claim provides the desired contradiction via Proposition~\ref{prop:non-escape}. We verify the claim in several cases:

\begin{enumerate}[label=\textbf{Case \arabic*:}]
\item Suppose $t\notin\bigcup_j J_j$. Then $t\notin K$, so either $\haus^n(\partial\Omega_t\cap M)<W_g(U)-1/i$ or $\haus^n(\partial\Omega_t\cap N_r(U))<\varepsilon(U)/2$. Since $\partial\Omega_t'=\partial\Omega_t$, the claim holds.

\item Suppose $t\in\bigcup_j J_j$ but $t\notin K$. If $\haus^n(\partial\Omega_t\cap M)<W_g(U)-1/i$, then since $t$ is contained in at most two intervals,
\[
\haus^n(\partial\Omega_t'\cap M)\le \haus^n(\partial\Omega_t\cap M)+2\cdot\frac{1}{4i}<W_g(U).
\]
If $\haus^n(\partial\Omega_t\cap N_r(U))<\varepsilon(U)/2$, then $\{\Omega_t'\}$ only differs from $\{\Omega_t\}$ inside $U_j\cup U_{j+1}$ (for at most two indices). If $U_j\cup U_{j+1}$ is disjoint from $U$, then $\haus^n(\partial\Omega_t'\cap U)=\haus^n(\partial\Omega_t\cap U)<\varepsilon(U)/2$. If $U_j$ intersects $U$, then by definition of $\mathcal{A}(r,U)$ we have $U_j\subset N_r(U)$, so $\haus^n(\partial\Omega_t\cap U_j)<\varepsilon(U)/2$. Thus
\[
\haus^n(\partial\Omega_t'\cap N_r(U))\le \haus^n(\partial\Omega_t\cap N_r(U))+2\cdot\frac{1}{4i}<\varepsilon(U)/2+\frac{1}{2i}<\varepsilon(U).
\]

\item Suppose $t\in K$. Since $\{(b_j+\delta,a_j-\delta)\}$ covers $K$ and each point is in at most two intervals,
\[
\haus^n(\partial\Omega_t'\cap M)\le \haus^n(\partial\Omega_t\cap M)+\frac{1}{4i}-\frac{1}{2i}\le W_g(U)-\frac{1}{8i}.
\]
\end{enumerate}
This proves the claim and completes the proof.
\end{proof}

\section{Limit varifolds and proof of the main theorem}\label{sec:main-proof}
Let $\{\Gamma_i\}$ be the sequence given by Proposition~\ref{prop:AM-selection}, and extract a subsequence (not relabelled) such that the induced varifolds converge to $V$.  Property (1) guarantees that $\{\Gamma_i\}$ is almost minimizing in small relative annuli, while property (3) prevents the sweepout from escaping the good set.  The free boundary regularity theory of Li--Zhou~\cite{LZ17} thus implies:

\begin{proposition}\label{prop:regular-limit}
There exist finitely many connected, smooth, almost properly embedded free boundary minimal hypersurfaces $(\Sigma_j,\partial\Sigma_j)\subset (M,\partial M)$ and positive integers $m_j$ such that
\[
V=\sum_{j=1}^N m_j|\Sigma_j|.
\]
\end{proposition}

Set $\Sigma:=\bigcup_{j=1}^N \Sigma_j$.  By property (3) in Proposition~\ref{prop:AM-selection}, $\haus^n(\Sigma\cap N_\delta(U))\ge \varepsilon(U)/2$ for every fixed $\delta>0$, while
\[
\haus^n(\Sigma\cap M)\le W_\partial(U)+\haus^n(\partial_{\mathrm{rel}}U).
\]

\begin{theorem}\label{thm:main-sup}
Let $U\subset M$ be a good set.  For every $\delta>0$ there exists a smooth free boundary minimal hypersurface $(\Sigma,\partial\Sigma)\subset (M,\partial M)$ with $\partial\Sigma\subset\partial M$ such that
\begin{align*}
    \haus^n(\Sigma\cap M) &\le W_\partial(U)+\haus^n(\partial_{\mathrm{rel}}U),\\
    \haus^n(\Sigma\cap N_\delta(U)) &\ge \frac{\varepsilon(U)}{2}.
\end{align*}
\end{theorem}

Combining Lemma~\ref{lem:good-set-sublinear}, Proposition~\ref{prop:non-escape}, Proposition~\ref{prop:regular-limit}, and Theorem~\ref{thm:main-sup}, we obtain the free boundary minimal hypersurface.  Since the sweepout represents a nontrivial class in $H_n(M,\partial M;\mathbb{Z}_2)$, every sweepout slice is a relative cycle with $\partial\Gamma_t\subset\partial M$ (condition (sw0)), and the min-max limit $\Sigma$ must satisfy $\partial\Sigma\neq\varnothing$.  Indeed, if $\partial\Sigma=\varnothing$ then $\Sigma$ would be a closed cycle in $M$, which represents the trivial class in $H_n(M,\partial M;\mathbb{Z}_2)$, contradicting the nontriviality of the sweepout class.  The Li--Zhou regularity theory~\cite{LZ17} then produces $\Sigma$ with $\partial\Sigma\subset\partial M$ and $\Sigma$ meeting $\partial M$ orthogonally.  This completes the proof of the main existence result.
